\documentclass[10pt,a4paper]{amsart}
\usepackage[margin=19mm]{geometry}
\usepackage{amsmath,amssymb,mathtools}
\usepackage[scr=rsfs,cal=euler]{mathalfa}
\usepackage{xfrac}
\usepackage[unicode,hidelinks,bookmarks=false]{hyperref}
\newcommand{\ie}{i.e.\ }
\newcommand{\cf}{cf.\ }
\newcommand{\resp}{respectively\ }
\newcommand{\Subsection}{\subsection}
\providecommand{\nobreakdash}{\nobreak\hbox{-}}
\setlength{\emergencystretch}{2em}
\multlinegap0pt
\usepackage{bm}
\usepackage{bbm}
\usepackage{dsfont}
\usepackage{xspace}
\usepackage{cancel}
\usepackage{mathtools}
\usepackage{amsthm}
\makeatletter
\@ifundefined{theorem}{\newtheorem{theorem}{Theorem}}{}
\@ifundefined{lemma}{\newtheorem{lemma}[theorem]{Lemma}}{}
\makeatother
\title[Residual Finiteness of $A\_{2,3,2n}$ Triangle Artin Groups]{Residual Finiteness of $A\_{2,3,2n}$ Triangle Artin Groups}
\author{Greyson Meyer}
\date{Source: arXiv:2412.07063v2 (CC BY 4.0)}
\begin{document}
\maketitle

\noindent\emph{Source and licence.} Residual Finiteness of $A_{2,3,2n}$ Triangle Artin Groups. Version arXiv:2412.07063v2 (CC BY 4.0), distributed under the Creative Commons Attribution 4.0 International licence, as recorded in the arXiv licence metadata for this exact version and shown on the abstract page https://arxiv.org/abs/2412.07063v2. Full rights notice: licenses/arxiv-2412.07063-CC-BY-4.0.txt.\par\smallskip

\noindent\textit{Editorial scope.} Counted unit: the Lemma whose source label is \texttt{None} in the article (source0.tex lines 294--296, source label \texttt{None}), delivered together with the article's own complete original proof of it (source0.tex lines 297--336, one \texttt{proof} environment of 40 lines). No mathematics is written, repaired or re-derived: statement and proof are the article's own text line for line. Required context retained uncounted: none. Every definition, display and label of those lines is retained unchanged. Omitted from this package and not redistributed: 1--293, 337--843. Nothing inside a retained excerpt is omitted or altered: every retained body is delivered exactly as the article's lines are written, so no omission receipt of any kind accompanies this package. Citations: the retained excerpts carry no citation command at all, so no bibliography entry of the article is transcribed and no reference is redistributed. Reference adaptation: none was needed and none was made; every reference inside the delivered unit resolves inside it, so no unresolved reference or citation is produced. Printed slips kept literal: no printed word, symbol, hypothesis or number of the counted statement or of its proof was altered. Layout and class: the article's own class is \texttt{nyjm} with packages including hyperref, graphicx, amssymb, upref, amsthm, xcolor, multirow, mathtools, tikz-cd; the wrapper is the lane's pinned \texttt{amsart} editorial wrapper at 10pt on A4, which restates the theorem-like environments the retained text needs and adds the article's own macro definitions that the retained text needs (none), with extra wrapper packages (bm, bbm, dsfont, xspace, cancel, mathtools). The delivered numbering is the unit's own. Assets: no retained span contains a figure environment, a graphics-inclusion command, a PDF-page inclusion, a table or a TikZ picture; no image file is copied into this package, the package declares no asset dependency and no image file is redistributed with it. Licence: version arXiv:2412.07063v2 of this article is distributed under the Creative Commons Attribution 4.0 International licence, as recorded in the arXiv licence metadata for this exact version and shown on the abstract page https://arxiv.org/abs/2412.07063v2. A full rights notice is delivered at licenses/arxiv-2412.07063-CC-BY-4.0.txt. Characters: every retained excerpt is pure ASCII, so the delivered engine, XeLaTeX with its default Latin Modern text font, reports no missing character.
\begin{lemma}
    For all $g\in A_{2,3,2n}$, $\pi_1(U)\cap\pi_1(\sigma(W))^g$ is contained in a $\pi_1(U)$-conjugacy class of subgroups of $\pi_1(U)$ represented by a labeled graph in $S$. Likewise, $\pi_1(V)\cap\pi_1(\sigma(W))^g$ is also contained in a $\pi_1(U)$-conjugacy class of subgroups of $\pi_1(U)$ represented by a labeled graph in $S$.
\end{lemma}

\begin{proof}
    Lemma 3.4 allows us to decompose every $g\in A_{2,3,2n}$ as \[g=g_1vg_2v...g_{m-1}vg_m\]
    where $m\geq1$,  $g_i\in\pi_1(U)$ and $v\in\pi_1(V)$ is the fixed nontrivial coset representative coming from the quotient $\pi_1(V)/\pi_1(W)\cong\mathbb{Z}/2\mathbb{Z}$. Define $\ell(g)$ to be the number of nontrivial elements in the decomposition of $g$ as above. We will proceed via induction on $\ell(g)$.

    Base case: Consider an element $g\in A_{2,3,2n}$ where $\ell(g)=1$. Then $g=v$ or $g$ is some nontrivial element in $\pi_1(U)$.
    
    First assume $g=v$. Since $\pi_1(W)^v=\pi_1(W)=\pi_1(\sigma(W))$, we have $\pi_1(U)\cap\pi_1(\sigma(W))^g=\pi_1(V)\cap\pi_1(\sigma(W))^g=\pi_1(\sigma(W))$.

    Now assume $g\in\pi_1(U)$. Then $\pi_1(\sigma(W))^g<\pi_1(U)$, making $\pi_1(U)\cap\pi_1(\sigma(W))^g=\pi_1(\sigma(W))^g$. The $\pi_1(U)$-conjugacy class of this subgroup is represented by $\sigma(W)$ with its basepoint shifted according to $g$. The labeled graph $\sigma(W)\in S$ by construction.

    In the Bass-Serre tree $T$ associated with the splitting of $A_{2,3,2n}$ arising from Theorem 2.4, we see that the vertex stabilized by $\pi_1(U)$ is between the edge stabilized by $\pi_1(\sigma(W))^g$ and the vertex stabilized by $\pi_1(V)$. Therefore $\pi_1(U)$ is also stabilized by $\pi_1(V)\cap\pi_1(\sigma(W))^g$. So we can rewrite $\pi_1(V)\cap\pi_1(\sigma(W))^g=\pi_1(V)\cap\pi_1(U)\cap\pi_1(\sigma(W))^g=\pi_1(\sigma(W))\cap\pi_1(\sigma(W))^g$. Therefore, by Lemma 3.2, the labeled graph associated with the $\pi_1(U)$-conjugacy class of this subgroup is a connected component of $\sigma(W)\otimes_U\sigma(W)$, all of which are contained in $S$ since $S$ is closed under fiber product.

    Induction hypothesis: Assume that the lemma holds for all $h\in A_{2,3,2n}$ with $\ell(h)<k$ for some $k>1$.
    
    Induction step: Consider an element $g\in A_{2,3,2n}$ with $\ell(g)=k$.
    
    Case 1:
Let $g=g'v$ where the decomposition of $g'$ begins with $v$ and $\ell(g')=k-1$. Consider the Bass-Serre tree $T$ induced by the splitting of $A_{2,3,2n}$. The subgroup $\pi_1(U)\cap\pi_1(\sigma(W))^g$ is the stabilizer of the path $p$ in $T$ between the vertex in $T$ stabilized by $\pi_1(U)$ and the edge in $T$ stabilized by $\pi_1(\sigma(W))^g$. Since $g$ begins with $v$, this makes the vertex stabilized by $\pi_1(U)^v$ on the path $p$. Therefore $\pi_1(U)\cap\pi_1(\sigma(W))^g=\pi_1(U)\cap(\pi_1(U)\cap\pi_1(\sigma(W))^{g'})^v$. By the induction hypothesis, $\pi_1(U)\cap\pi_1(\sigma(W))^{g'}$ is contained in a $\pi_1(U)$-conjugacy class represented by a labeled graph $H^1_U\in S$. The $v$-conjugate of a subgroup contained in the $\pi_1(U)$-conjugacy class represented by $H^1_U$ is contained in the $\pi_1(U)$-conjugacy class represented by $\beta(H^1_U)$ by the definition of $\beta$. The graph $\beta(H^1_U)\in S$ since $S$ is closed under $\beta$. 


Similarly, $\pi_1(V)\cap\pi_1(\sigma(W))^g=(\pi_1(V)\cap\pi_1(\sigma(W))^{g'})^v$ since $v\in\pi_1(V)$. By the induction hypothesis, $\pi_1(V)\cap\pi_1(\sigma(W))^{g'}$ is contained in a $\pi_1(U)$-conjugacy class represented by a labeled graph $H^1_V\in S$. Therefore $\pi_1(V)\cap\pi_1(\sigma(W))^g$ is represented by $\beta(H^1_V)\in S$.

Case 2: Now assume that $g=g'v$ and the decomposition of $g'$ begins with some nontrivial $g_1\in\pi_1(U)$. Since $k>1$, this forces $g$ to begin with $g_1v$. The vertex in $T$ stabilized by $\pi_1(U)^{g_1v}=\pi_1(U)^v$ is on the path in $T$ between the vertex in $T$ stabilized by $\pi_1(U)$ and the edge in $T$ stabilized by $\pi_1(\sigma(W))^g$. Therefore $\pi_1(U)\cap\pi_1(\sigma(W))^g=\pi_1(U)\cap\pi_1(U)^{g_1v}\cap\pi_1(\sigma(W))^g=\pi_1(U)\cap(\pi_1(U)^{g_1}\cap\pi_1(\sigma(W))^{g'})^v=\pi_1(U)\cap(\pi_1(U)\cap\pi_1(\sigma(W))^{g'})^v$. By the induction hypothesis, $\pi_1(U)\cap\pi_1(\sigma(W))^{g'}$ is contained in the $\pi_1(U)$-conjugacy class represented by a labeled graph $H^2_U\in S$. So $\pi_1(U)\cap\pi_1(\sigma(W))^g=\pi_1(U)\cap(\pi_1(U)\cap\pi_1(\sigma(W))^{g'})^v$ is represented by $\beta(H^2_U)\in S$.


Similarly, $\pi_1(V)\cap\pi_1(\sigma(W))^g=(\pi_1(V)\cap\pi_1(\sigma(W))^{g'})^v$. By the induction hypothesis we know that $\pi_1(V)\cap\pi_1(\sigma(W))^{g'}$ is contained in the $\pi_1(U)$-conjugacy class represented by a labeled graph $H^2_V\in S$. Therefore $\pi_1(V)\cap\pi_1(\sigma(W))^g$ is represented by $\beta(H^2_V)\in S$.


Case 3: Assume $g=g'g_m$, $g_m\neq1$, the decomposition of $g'$ begins with $v$ and $\ell(g')=k-1$. Then $\pi_1(U)\cap\pi_1(\sigma(W))^g=(\pi_1(U)\cap\pi_1(\sigma(W))^{g'})^{g_m}$ since $g_m\in\pi_1(U)$. By the induction hypothesis we know that $\pi_1(U)\cap\pi_1(\sigma(W))^{g'}$ is contained in the $\pi_1(U)$-conjugacy class represented by a labeled graph $H^3_U\in S$. Therefore $\pi_1(U)\cap\pi_1(\sigma(W))^g$ is also represented by $H^3_U$.


Define $g''$ to be such that $g=vg''$ with $\ell(g'')=k-1$ and $g''$ ends with $g_m$. Then $\pi_1(V)\cap\pi_1(\sigma(W))^g=\pi_1(V)\cap((\pi_1(V)\cap\pi_1(\sigma(W)))^v)^{g''}=\pi_1(V)\cap(\pi_1(V)\cap\pi_1(\sigma(W))^v)^{g''}$ since $v\in\pi_1(V)$. The subgroup $\pi_1(V)\cap\pi_1(\sigma(W))^v=\pi_1(\sigma(W))$ since $\pi_1(\sigma(W))=\pi_1(W)\unlhd\pi_1(V)$, so $\pi_1(V)\cap\pi_1(\sigma(W))^g=\pi_1(V)\cap\pi_1(\sigma(W))^{g''}$ which is contained in the $\pi_1(U)$-conjugacy class represented by a labeled graph in $S$ by the induction hypothesis.

Case 4: Now assume that $g=g'g_m$ and the decomposition of $g'$ begins with some nontrivial $g_1\in\pi_1(U)$. Then $\pi_1(U)\cap\pi_1(\sigma(W))^g=(\pi_1(U)\cap\pi_1(\sigma(W))^{g'})^{g_m}$ since $g_m\in\pi_1(U)$. By the induction hypothesis, $\pi_1(U)\cap\pi_1(\sigma(W))^{g'}$ is contained in a $\pi_1(U)$-conjugacy class represented by a labeled graph $H^4_U\in S$. Therefore $\pi_1(U)\cap\pi_1(W)^g$ is also represented by $H^4_U\in S$.


Since $g'$ begins with the nontrivial $g_1\in\pi_1(U)$, the path in $T$ stabilized by $\pi_1(V)\cap\pi_1(\sigma(W))^g$ passes through the vertex in $T$ stabilized by $\pi_1(U)$. Therefore $\pi_1(V)\cap\pi_1(\sigma(W))^g=\pi_1(V)\cap\pi_1(U)\cap\pi_1(\sigma(W))^g$. The previous step tells us that $\pi_1(U)\cap\pi_1(\sigma(W))^g$ is represented by the labeled graph $H^4_U\in S$, making $\pi_1(U)\cap\pi_1(\sigma(W))^g\leq\pi_1(U)$. Since $\pi_1(U)\cap\pi_1(V)=\pi_1(\sigma(W))$, $\pi_1(V)\cap\pi_1(\sigma(W))^g=\pi_1(\sigma(W))\cap\pi_1(\sigma(W))^g$, forcing $\pi_1(V)\cap\pi_1(\sigma(W))^g$ to be contained in the $\pi_1(U)$-conjugacy class represented by a connected component of $\sigma(W)\otimes_UH_U^4\in S$.

Therefore $\pi_1(U)\cap\pi_1(\sigma(W))^g$ and $\pi_1(V)\cap\pi_1(\sigma(W))^g$ are each contained in a $\pi_1(U)$-conjugacy class represented by a labeled graph in $S$ for all $g\in A_{2,2,3n}$.
\end{proof}

\end{document}
